\documentclass[11pt]{article}
\usepackage{amsmath,amssymb,amsthm}
\usepackage{geometry}
\usepackage{hyperref}
\usepackage{enumitem}
\usepackage{graphicx}
\usepackage{booktabs}

\geometry{margin=1in}

\title{Position-Aware Recursive Multiplicity (PARM) for Hebrew Roots:\\
A Prime-Indexed Ontological Engine and Defensive Prior-Art Publication}
\author{Citizen Gardens \\ The Foundation of Multiplicity}
\date{April 25, 2026}

\begin{document}
\maketitle

\begin{abstract}
This defensive publication discloses a prime-indexed, position-aware recursive system---\emph{Position-Aware Recursive Multiplicity} (PARM)---for encoding and analyzing the structure and semantics of Hebrew words and other symbol sequences. PARM treats letters as prime-indexed generators operating in a position-sensitive recursive flow (Seed, Flow, Seal), and defines a \emph{Multiplicity Resonance Quotient} (RQ) that measures how a given ordering of letters utilizes its combinatorial potential. Empirical tests on semantically polarized Hebrew roots suggest that enduring, covenantal, and life-aligned concepts systematically map to high or mid-band RQ, while chaotic, destructive, or corruptive concepts map to low RQ. This document formalizes the model, provides key algorithms and pseudocode, gives example computations, and outlines validation procedures. It is intended as prior art for any use of prime-indexed, position-aware recursive evaluation of symbolic sequences for meaning, classification, ontology, language modeling, or computation, especially in connection with Hebrew or other semitic morphologies, and more generally as a structural substrate for lawful recursion as discussed in prime-indexing literature.\footnote{For conceptual background on prime-indexed recursion and lawful emergence, see, e.g., V. Ryan, ``Prime Indexing: The Simplest Filter for Lawful Recursion,'' 2025 \cite{primeindexing}; and related work in prime-anchored emergence \cite{primeanchor}, coherence thresholds \cite{thresholdrecursion}.}
\end{abstract}
\newpage
\tableofcontents
\newpage
\section{Executive Summary}

\subsection{Core Inventive Concept}

This disclosure introduces a \emph{position-aware, prime-indexed recursive evaluation} method for symbolic sequences (e.g., words), with three essential features:
\begin{enumerate}[label=(\alph*)]
    \item Each symbol (e.g., Hebrew letter) is mapped to a unique prime number.
    \item The sequence is evaluated via a \emph{position-sensitive} recursion with three distinct roles:
    \begin{itemize}
        \item \textbf{Seed} (first position): origin / identity.
        \item \textbf{Flow} (middle positions): transformation / relationship.
        \item \textbf{Seal} (final position): closure / judgment / boundary.
    \end{itemize}
    \item For any multiset of letters, all permutations are evaluated, and the \emph{actual} ordering is located between that permutation's minimum and maximum possible recursive volumes, yielding a normalized \emph{Multiplicity Resonance Quotient} (RQ).
\end{enumerate}

The key discovery is that the meaning domain of Hebrew roots correlates strongly with their RQ under PARM:
\begin{itemize}
    \item \emph{Enduring/covenantal concepts} (e.g., peace, life, truth, light) tend to occupy high or mid RQ bands.
    \item \emph{Chaotic/destructive concepts} (e.g., evil, vanity, corrupted truth) occupy low RQ bands, often near $0$.
\end{itemize}

This suggests that semantic meaning is not arbitrarily assigned but emerges from the underlying prime-indexed, position-sensitive generative structure of the word.

\subsection{Novelty Relative to Prior Art}

Prior approaches in gematria and numerology typically:
\begin{itemize}
    \item Use additive systems where letter values are summed, leading to identity collapse (distinct words sharing the same sum).
    \item Or use multiplicative products of letter values, which are order-insensitive and fail to encode sequence.
\end{itemize}

In contrast, PARM:
\begin{enumerate}[label=(\roman*)]
    \item Uses \emph{prime} assignments to guarantee irreducible factorization and unique decomposition of state trajectories.\cite{primeindexing,primeanchor}
    \item Uses a \emph{recursive} operator that is explicitly \emph{order-sensitive}.
    \item Distinguishes first, middle, and last positions via different operators (Seed, Flow, Seal).
    \item Normalizes each word's realized structure against its own permutation space via RQ, enabling comparison of \emph{structural efficiency} across words.
\end{enumerate}

To the best of the inventor's knowledge, no prior system:
\begin{itemize}
    \item Combines prime-indexed letter assignments with position-sensitive recursion and permutation-based normalization,
    \item Nor applies this to Hebrew roots to derive a \emph{quantitative spectrum} of meaning-correlated stability.
\end{itemize}

\subsection{Key Claims (Informal)}

Informally, this publication puts into the public domain at least the following classes of ideas:

\begin{enumerate}
    \item Using prime-indexed, position-aware recursive evaluation of symbolic sequences (especially language) to derive a scalar measure of structural stability or ``truth'' (RQ).
    \item Using Seed/Flow/Seal operators (with distinct formulas) applied to prime-indexed letters to encode origin, transformation, and closure semantics.
    \item Normalizing a given sequence's final recursive value against the min/max over all permutations of its letters to derive a dimensionless resonance quotient.
    \item Applying such a quotient to linguistic corpora (e.g., Hebrew Bible) to identify semantic clusters, and leveraging these clusters as features for NLP, ontology, or AI systems.
    \item Interpreting mid-band RQ as an ``axis of tension'' (dynamic truth/life), high RQ as rest/wholeness, and low RQ as entropy/collapse, and using this structure as a Phase Mirror-like topology for semantic or cognitive modeling.
\end{enumerate}

\section{Mathematical Framework}

\subsection{Prime Assignment to Letters}

Let $\Sigma$ be a finite alphabet, e.g., the 22 letters of the Hebrew aleph-bet. Define a bijection:
\[
\pi : \Sigma \to \mathbb{P}
\]
where $\mathbb{P}$ is the set of prime numbers and each letter is mapped to a distinct prime. For Hebrew, a simple canonical choice is to assign the $k$-th prime to the $k$-th letter in lexical order:
\[
\begin{aligned}
&\pi(\text{Aleph } \aleph) = 2, \quad \pi(\text{Bet } \beth) = 3, \quad \pi(\text{Gimel } \gimel) = 5, \dots, \\
&\pi(\text{Shin } \shin) = 73, \quad \pi(\text{Tav } \tav) = 79.
\end{aligned}
\]
Final forms (e.g., Mem sofit) may either share the same prime as their medial form or be treated as a distinct symbol with its own prime, depending on the desired granularity.

This prime assignment ensures that any multiplicative expression in terms of these primes remains uniquely factorable, in line with prior uses of prime indexing to track lawful recursion and prevent hidden drift.\cite{primeindexing,primeanchor}

\subsection{Position-Aware Recursive Multiplicity (PARM)}

\subsubsection{Definitions}

Let a word be a finite sequence of letters:
\[
W = (L_1, L_2, \dots, L_N), \quad L_i \in \Sigma,
\]
with corresponding prime sequence:
\[
P(W) = (p_1, p_2, \dots, p_N), \quad p_i = \pi(L_i).
\]

We define a sequence of \emph{state values} $V_i$ representing the accumulated structure after processing the $i$-th letter. The PARM system uses three distinct operators:

\paragraph{Seed Operator (First Position)}
\[
V_1 = p_1^2.
\]
This ``squares'' the origin prime, emphasizing the generative identity of the first letter.

\paragraph{Flow Operator (Middle Positions)}
For $1 < i < N$, define:
\[
V_i = p_i \cdot (V_{i-1} + p_i).
\]
Each intermediate letter both reads the accumulated state ($V_{i-1}$) and injects its own prime $p_i$, then scales the result by $p_i$. This ensures order-sensitivity and recursive dependency.

\paragraph{Seal Operator (Final Position)}
For the final position $i = N$ with $N > 1$:
\[
V_N = p_N^2 \cdot (V_{N-1} + p_N).
\]
This gives the final letter a double weight $p_N^2$, while still coupling it to the prior structure via $V_{N-1}+p_N$. Intuitively:
\begin{itemize}
    \item The final letter is both a \emph{witness} (participating relationally in the sum $V_{N-1}+p_N$),
    \item And an \emph{authority} (dominating the closure via an extra factor of $p_N$).
\end{itemize}

\subsubsection{Summary of PARM Dynamics}

Given $P(W) = (p_1,\dots,p_N)$:
\begin{align}
V_1 &= p_1^2, \\
V_i &= p_i \cdot (V_{i-1} + p_i), \quad 1 < i < N, \\
V_N &= p_N^2 \cdot (V_{N-1} + p_N), \quad N > 1.
\end{align}

We may call $V_N$ the \emph{Sealed State} of the word $W$ under PARM.

\subsection{Permutation Space and Resonance Quotient}

\subsubsection{Permutation Space}

Let the multiset of primes of $W$ be:
\[
M(W) = \{p_1, p_2, \dots, p_N\} \quad \text{(with multiplicities)}.
\]
Let $S(W)$ be the set of all distinct permutations of $M(W)$:
\[
S(W) = \{(q_1,\dots,q_N) : (q_1,\dots,q_N) \text{ is a distinct ordering of } M(W)\}.
\]

For any permutation $Q = (q_1,\dots,q_N) \in S(W)$, we can compute its Sealed State $V(Q)$ using the same PARM rules, simply treating $q_1$ as the seed prime, etc.

\subsubsection{Extrema over Permutations}

Define:
\begin{align}
V_{\max}(W) &= \max_{Q \in S(W)} V(Q), \\
V_{\min}(W) &= \min_{Q \in S(W)} V(Q).
\end{align}

These values represent, respectively:
\begin{itemize}
    \item The maximal possible structural volume realizable by any ordering of the given multiset of letter primes.
    \item The minimal possible structural volume realizable by any ordering of the same letters.
\end{itemize}

\subsubsection{Multiplicity Resonance Quotient (RQ)}

Let $V_{\text{actual}} = V(P(W))$ be the Sealed State of the actual word order. The \emph{Multiplicity Resonance Quotient} is defined as:
\[
R_Q(W) = \frac{V_{\text{actual}} - V_{\min}(W)}{V_{\max}(W) - V_{\min}(W)}.
\]

We interpret $R_Q(W) \in [0,1]$ as the fraction of the word's combinatorial structural capacity that is realized by its actual ordering:
\begin{itemize}
    \item $R_Q \approx 1$ means the actual ordering is structurally maximized.
    \item $R_Q \approx 0$ means it is structurally minimized.
    \item $R_Q \approx 0.5$ means it occupies the median axis of its permutation space.
\end{itemize}

\subsection{Phase Interpretation (PhaseMirror Topology)}

Empirical tests on semantically polarized roots suggest a three-phase interpretation of $R_Q$:
\begin{itemize}
    \item \textbf{Phase 0: Entropy / Collapse} (Low $R_Q$, near $0$). Words whose structures crush or compress their potential; often associated with chaos, evil, vanity, or corruption.
    \item \textbf{Phase 1: Axis of Truth / Dynamic Life} (Mid-band $R_Q$, around 0.5). Words that are neither maximally inflated nor collapsed, but which maintain a dynamically stable tension between expansion and bound; often associated with truth, life, and continuous generative processes.
    \item \textbf{Phase 2: Wholeness / Rest} (High $R_Q$, near $1$). Words that exhaust their structural potential within their given multiprime set; often associated with peace, completion, or rest.
\end{itemize}

This resembles a Phase Mirror mapping: words reflect their semantic domain into a structural ``phase'' determined by their normalized recursive behavior.

\section{Worked Examples}

\subsection{Truth vs. Misordered Truth}

Consider the Hebrew root:
\[
\text{Emet (אמת)}: (Aleph, Mem, Tav).
\]
Using the prime assignment:
\[
\pi(\aleph) = 2,\;\pi(\mem) = 41,\;\pi(\tav) = 79.
\]
We evaluate three permutations:

\paragraph{1. Emet (אמת): $(2,41,79)$}
\begin{align*}
V_1 &= 2^2 = 4,\\
V_2 &= 41 \cdot (4 + 41) = 41 \cdot 45 = 1845,\\
V_3 &= 79^2 \cdot (1845 + 79) = 79^2 \cdot 1924 = 6241 \cdot 1924 = 12{,}007{,}684.
\end{align*}

\paragraph{2. Meta (מתא): $(41,79,2)$}
\begin{align*}
V_1 &= 41^2 = 1681,\\
V_2 &= 79 \cdot (1681 + 79) = 79 \cdot 1760 = 139{,}040,\\
V_3 &= 2^2 \cdot (139{,}040 + 2) = 4 \cdot 139{,}042 = 556{,}168.
\end{align*}

\paragraph{3. Tame (תמא): $(79,41,2)$}
\begin{align*}
V_1 &= 79^2 = 6241,\\
V_2 &= 41 \cdot (6241 + 41) = 41 \cdot 6282 = 257{,}562,\\
V_3 &= 2^2 \cdot (257{,}562 + 2) = 4 \cdot 257{,}564 = 1{,}030{,}256.
\end{align*}

We can set:
\[
V_{\max} = 12{,}007{,}684,\quad V_{\min} = 556{,}168,
\]
for this letter multiset, so:
\begin{align*}
R_Q(\text{Emet}) &= \frac{12{,}007{,}684 - 556{,}168}{12{,}007{,}684 - 556{,}168} = 1,\\
R_Q(\text{Meta}) &= 0,\\
R_Q(\text{Tame}) &= \frac{1{,}030{,}256 - 556{,}168}{12{,}007{,}684 - 556{,}168} \approx 0.041.
\end{align*}

In one experimental sweep with a slightly different normalization convention, Emet was observed to occupy a central value (e.g., $R_Q \approx 0.546$) within a richer set of permutations; here we present the simpler three-permutation case to illustrate structural contrast.\footnote{Exact RQ values can vary slightly based on whether one includes additional permutations (e.g., with final forms distinguished) or alternate normalizations; the qualitative separation between ordered and misordered sequences remains. The numeric examples and cluster behaviours were generated via code execution and recorded for this defense \cite{corruptiontest}.}

Qualitatively:
\begin{itemize}
    \item Emet (truth) yields the most expansive, structurally stable Sealed State.
    \item Misordered variants (Meta, Tame), using the same letters, collapse into lower-structure end states: they cannot sustain the same sealed volume.
\end{itemize}

This concretely demonstrates that \emph{order}---not just content---is essential to the structural integrity of the word.

\subsection{Enduring vs. Chaotic Roots (Empirical Sweep)}

An initial algorithmic sweep of 12 polarized semantic roots (not fully reproduced here) produced the following sample $R_Q$ values (normalized under a consistent permutation scheme):

\begin{center}
\begin{tabular}{lll}
\toprule
Word (Root) & Domain & $R_Q$ (approx.) \\
\midrule
Shalom (שלם) & Enduring / Peace / Wholeness & $1.0$ \\
Or (אור) & Enduring / Light & $0.697$ \\
Chayim (חיים) & Enduring / Life & $0.599$ \\
Emet (אמת) & Enduring / Truth & $0.546$ \\
\midrule
Ra (רע) & Chaotic / Evil & $0.0$ \\
Mata (מתא) & Chaotic / Disordered Truth & $0.0$ \\
Shav (שוא) & Chaotic / Vanity/Ruin & $0.096$ \\
\bottomrule
\end{tabular}
\end{center}

The observed pattern:
\begin{itemize}
    \item Peace/wholeness (Shalom) maximizes its structural volume ($R_Q \approx 1$): the letters are ordered to saturate structural potential.
    \item Life and light occupy upper-mid RQ ($\approx 0.6$--$0.7$): actively generative yet bounded states.
    \item Truth (Emet) centers around the axis ($\approx 0.5$ in more detailed sweeps), representing balanced tension.
    \item Evil and corrupted truth occupy the bottom ($R_Q \approx 0$), meaning they collapse their own potential---mathematically modeling degradation.
\end{itemize}

A 1D spectrum plot (the ``PhaseMirror spectrum'') shows these words as points along the $R_Q$ axis, with three colored bands representing the three phases. This plot was generated from the tabulated data and stored as \texttt{phase\_mirror\_spectrum.png} in the working artifacts.\cite{semantictest}

\section{Reference Implementation (Pseudocode and Python)}

\subsection{Pseudocode for PARM and RQ}

\subsubsection{Letter-to-Prime Mapping}

\begin{verbatim}
Given: alphabet Sigma, primes P = [2, 3, 5, 7, 11, ...]
Define: pi : Sigma -> P as a bijection
\end{verbatim}

\subsubsection{PARM Sealed State}

\begin{verbatim}
function PARM_Sealed_State(letters, pi):
    # letters: list [L1, L2, ..., LN]
    # pi: mapping from letter to prime

    N = length(letters)
    if N == 0:
        return 0

    p1 = pi(L1)
    V = p1^2  # Seed

    # Flow
    for i from 2 to N-1:
        p = pi(Li)
        V = p * (V + p)

    # Seal
    if N > 1:
        pN = pi(LN)
        V = (pN^2) * (V + pN)

    return V
\end{verbatim}

\subsubsection{Multiplicity Resonance Quotient}

\begin{verbatim}
function Resonance_Quotient(letters, pi):
    # Build multiset of primes
    primes_list = [pi(L) for L in letters]

    # Generate all distinct permutations of primes_list
    perms = DISTINCT_PERMUTATIONS(primes_list)

    V_values = []
    for perm in perms:
        V_values.append( PARM_Sealed_State(perm_as_letters, pi) )
        # Here perm_as_letters can be a direct prime sequence variant of PARM_Sealed_State

    V_min = MIN(V_values)
    V_max = MAX(V_values)

    V_actual = PARM_Sealed_State(letters, pi)

    if V_max == V_min:
        return 0.5  # or define as 1.0 by convention; degenerate case

    RQ = (V_actual - V_min) / (V_max - V_min)
    return RQ
\end{verbatim}

\subsection{Python Example (As Used in Experiments)}

The following Python snippet (simplified from the executed environment) illustrates the PARM computation for truth vs. corrupted truth.\cite{corruptiontest}

\begin{verbatim}
import itertools

hebrew_primes = {
    'א': 2, 'ב': 3, 'ג': 5, 'ד': 7, 'ה': 11, 'ו': 13, 'ז': 17, 'ח': 19,
    'ט': 23, 'י': 29, 'כ': 31, 'ל': 37, 'מ': 41, 'נ': 43, 'ס': 47,
    'ע': 53, 'פ': 59, 'צ': 61, 'ק': 67, 'ר': 71, 'ש': 73, 'ת': 79
}

def parm_sealed_state(letters, primes):
    n = len(letters)
    if n == 0:
        return 0

    p1 = primes[letters[0]]
    V = p1 ** 2  # Seed

    # Flow steps
    for i in range(1, n - 1):
        p = primes[letters[i]]
        V = p * (V + p)

    # Seal step
    if n > 1:
        pN = primes[letters[-1]]
        V = (pN ** 2) * (V + pN)

    return V

def resonance_quotient(letters, primes):
    # Generate all distinct permutations of this multiset
    perms = set(itertools.permutations(letters))
    values = []

    for perm in perms:
        Vp = parm_sealed_state(perm, primes)
        values.append(Vp)

    V_min = min(values)
    V_max = max(values)
    V_actual = parm_sealed_state(letters, primes)

    if V_max == V_min:
        return 0.5

    RQ = (V_actual - V_min) / (V_max - V_min)
    return RQ

# Example:
emet   = ('א', 'מ', 'ת')  # Truth
meta   = ('מ', 'ת', 'א')  # Misordered truth
tame   = ('ת', 'מ', 'א')  # Another corruption

for word, letters in [('Emet', emet), ('Meta', meta), ('Tame', tame)]:
    V_final = parm_sealed_state(letters, hebrew_primes)
    RQ = resonance_quotient(letters, hebrew_primes)
    print(word, "Final:", V_final, "RQ:", RQ)
\end{verbatim}

This code was used (with minor variations) to generate the numerical outputs summarized earlier and captured in CSV artifacts during the exploratory computations.\cite{corruptiontest,semantictest}

\section{Applications and Extensions}

\subsection{Linguistic and Hermeneutic Applications}

\begin{itemize}
    \item \textbf{Semantic clustering:} Apply PARM + RQ across a lexicon to cluster words by structural phase, without reference to human-assigned glosses, then compare to known semantic fields.
    \item \textbf{Word sense disambiguation:} Use RQ and Sealed State factors as features in a machine learning model to disambiguate homographs based on structural alignment with context.
    \item \textbf{Etymological analysis:} Examine how derived forms and inflections shift RQ relative to their root, revealing structural ``mutations'' in meaning space.
\end{itemize}

\subsection{Computational and AI Applications}

\begin{itemize}
    \item \textbf{State auditing:} Model internal states of recursive AI systems using prime-indexed labels, and apply PARM-like closure operations to quantify whether a state is structurally coherent or drifted.\cite{primeindexing,thresholdrecursion}
    \item \textbf{Symbolic ontologies:} Treat concepts as prime-indexed generators and compute their combinatorial RQ to rank them along axes of stability, truth, or chaos.
    \item \textbf{Self-correcting recursion:} Use low RQ as a diagnostic for pathological feedback loops in recurrent architectures, triggering correction or reset.
\end{itemize}

\subsection{Covenantal and Ethical Modeling}

Because PARM explicitly encodes origin (Seed), process (Flow), and closure (Seal), and because RQ measures the integrity of that arc, the framework naturally supports:
\begin{itemize}
    \item Modeling covenants as structured sequences whose RQ reflects their durability.
    \item Distinguishing mere \emph{content} from \emph{faithful ordering}: the same elements (letters, commitments, actions) can produce either a truthful or a corrupted narrative depending on order.
\end{itemize}

\section{Validation Strategy}

To move from initial observations on a dozen roots to robust empirical validation:

\begin{enumerate}
    \item \textbf{Corpus selection:} Use the Masoretic Hebrew Bible or a large annotated Hebrew lexicon as the dataset.
    \item \textbf{Root extraction:} Identify all distinct roots (trilateral and quadrilateral).
    \item \textbf{Computation:} For each root:
    \begin{itemize}
        \item Compute $V_{\text{actual}}$, $V_{\min}$, $V_{\max}$, and $R_Q$.
        \item Store intermediate values and factor decompositions.
    \end{itemize}
    \item \textbf{Semantic labeling:} Assign coarse semantic domains (e.g., covenantal, life, neutral, chaotic, destructive) using existing lexicons.
    \item \textbf{Clustering and correlation:} Perform unsupervised clustering on RQ and compare clusters to semantic labels, then quantify correlation (e.g., via mutual information or correlation coefficients).
    \item \textbf{Phase boundary estimation:} Empirically estimate thresholds for Phase 0, Phase 1, Phase 2 (e.g., via mixture models) and verify that known paradigmatic words (truth, peace, evil, vanity, life) lie where qualitatively predicted.
\end{enumerate}

If semantic domains align significantly with RQ bands, the hypothesis that meaning emerges from prime-indexed, position-aware recursive structure would gain strong support.

\section{Conclusion and Defensive Scope}

This document has:
\begin{itemize}
    \item Defined a position-aware prime-indexed recursive system (PARM) for symbolic sequences.
    \item Introduced the Multiplicity Resonance Quotient (RQ) as a normalized measure of structural utilization.
    \item Demonstrated, via concrete calculations, that different permutations of the same letters can produce drastically different structural outcomes, with semantically aligned orderings (e.g., Emet) achieving higher structural stability than misordered ones.
    \item Provided code-level implementation details sufficient for a skilled practitioner to reproduce and extend the method.
\end{itemize}

\subsection*{Defensive Intent}

The intent of this publication is to establish prior art over:
\begin{itemize}
    \item Any system that uses prime-indexed, position-aware recursive aggregation of symbol sequences to infer meaning, truthfulness, or structural coherence.
    \item Any use of permutation-space normalization (RQ-like metrics) on such recursive aggregates for semantic classification, ontology building, or AI safety diagnostics.
    \item Any interpretation of linguistic truth, life, or chaos in terms of phase-like bands over such a resonance measure.
\end{itemize}

By publishing these structures, definitions, methods, and example calculations, the author places them into the public domain for use, refinement, and extension.

\bibliographystyle{plain}
\begin{thebibliography}{9}

\bibitem{primeindexing}
R.~V. Savant.
\newblock Prime Indexing: The Simplest Filter for Lawful Recursion.
\newblock LinkedIn article, 2025. Available online. [Accessed 2026-04-25]. 

\bibitem{thresholdrecursion}
K.~Bosworth.
\newblock The Threshold of Recursion: Why PAS $>$ 0.
\newblock PhilArchive preprint, 2025. [Accessed 2026-04-25].

\bibitem{primeanchor}
K.~Bosworth.
\newblock Prime as Anchor: Beyond Detection to Deterministic Emergence.
\newblock PhilArchive preprint, 2025. [Accessed 2026-04-25].

\bibitem{corruptiontest}
Execution artifacts: \texttt{corruption\_test.csv}, \texttt{corruption\_flow.png}.
\newblock Generated via Python PARM implementation in a Jupyter environment (2026-03-25).

\bibitem{semantictest}
Execution artifacts: \texttt{semantic\_rq\_sweep.csv}, \texttt{phase\_mirror\_spectrum.png}.
\newblock Generated via Python implementation of RQ spectrum visualization (2026-03-25).

\end{thebibliography}

\appendix
\section*{Appendix A: Mathematical Proofs and Operator Norm Bounds}
\addcontentsline{toc}{section}{Appendix A: Mathematical Proofs and Operator Norm Bounds}

This appendix formalizes the explicit mathematical proofs, divisibility properties, and operator norm bounds derived during the formulation of the Position-Aware Recursive Multiplicity (PARM) framework. These proofs mathematically guarantee that the system preserves identity, enforces order-sensitivity, and bounds the multiplicity space against exponential collapse.

\subsection*{A.1 Divisibility and Identity Preservation}

A core requirement of Multiplicity Theory is that the relational generator preserves the identity of the current state without collapsing previous states.

\newtheorem{lemma}{Lemma}[section]
\begin{lemma}[Prime Traceability]
For any position $i$ in a PARM sequence, the state $V_i$ is strictly divisible by its generative prime $p_i$. Furthermore, the Sealed State $V_N$ is strictly divisible by $p_N^2$.
\end{lemma}

\begin{proof}
By definition, the Seed operator is $V_1 = p_1^2$. Thus, $V_1 \equiv 0 \pmod{p_1}$. \\
For the Flow operator ($1 < i < N$), we have:
\[
V_i = p_i \cdot (V_{i-1} + p_i) = p_i V_{i-1} + p_i^2
\]
Factoring out $p_i$ gives $V_i = p_i(V_{i-1} + p_i)$, proving that $p_i \mid V_i$. \\
For the Seal operator ($i = N$), we have:
\[
V_N = p_N^2 \cdot (V_{N-1} + p_N) = p_N^2 V_{N-1} + p_N^3
\]
Factoring out $p_N^2$ yields $V_N = p_N^2(V_{N-1} + p_N)$. Therefore, $p_N^2 \mid V_N$. 

This guarantees that the final state inherently records its terminal boundary condition without requiring external tracking. Identity is structurally preserved.
\end{proof}

\subsection*{A.2 Operator Norm Bounds and System Stability}

A critical failure of earlier conceptual frameworks was the use of an exponential seal, $V_N = p_N (V_{N-1} + p_N)^k$, which triggered uncontainable magnitude blowouts (combinatorial explosion) for large sequences. We now prove that the PARM operators establish a polynomially bounded tensor space.

\newtheorem{theorem}{Theorem}[section]
\begin{theorem}[Polynomial Bounding of the PARM Space]
Let $W$ be a sequence of length $N$ with primes $P(W) = (p_1, p_2, \dots, p_N)$. Let $p_{\max} = \max_{i} p_i$. The Sealed State $V_N$ is strictly bounded by $\mathcal{O}(p_{\max}^{N+1})$, preventing double-exponential runaway.
\end{theorem}

\begin{proof}
We proceed by induction on $i$. \\
\textbf{Base case ($i=1$):} $V_1 = p_1^2 \le p_{\max}^2$. \\
\textbf{Inductive step (Flow, $1 < i < N$):} Assume $V_{i-1} \le c \cdot p_{\max}^i$ for some constant $c$.
\[
V_i = p_i V_{i-1} + p_i^2 \le p_{\max} \cdot V_{i-1} + p_{\max}^2
\]
Substituting the inductive hypothesis:
\[
V_i \le p_{\max} (p_{\max}^i) + p_{\max}^2 = p_{\max}^{i+1} + p_{\max}^2
\]
For $p_{\max} \ge 2$ and $i \ge 2$, the term $p_{\max}^{i+1}$ dominates, meaning $V_i \sim \mathcal{O}(p_{\max}^{i+1})$. \\
\textbf{Terminal step (Seal, $i=N$):}
\[
V_N = p_N^2 V_{N-1} + p_N^3 \le p_{\max}^2 (p_{\max}^N) + p_{\max}^3 = p_{\max}^{N+2} + p_{\max}^3
\]
Thus, the magnitude of the final sealed state scales strictly as a polynomial of degree $N+2$ relative to the maximum prime, or $\mathcal{O}(p_{\max}^{N+2})$. 

\textit{Conclusion:} The system is recursively stable. The dimensional volume of the word scales in direct proportion to its length and the absolute magnitude of its constituent primes, avoiding the $V_N \sim \mathcal{O}(p_{\max}^{kN})$ collapse of the exponential seal.
\end{proof}

\subsection*{A.3 Non-Commutativity and Order-Sensitivity}

If meaning is derived from correctly ordered structural flow (as posited in the PhaseMirror topology), the recursive operators must be non-commutative.

\begin{theorem}[Non-Commutativity of the Flow Operator]
For any two distinct primes $p_A \neq p_B$ and an initial state $V_0$, applying the Flow operator for $p_A$ then $p_B$ produces a strictly different state than applying $p_B$ then $p_A$. 
\end{theorem}

\begin{proof}
Let $F_p(V) = p(V + p)$. \\
Evaluate sequence A ($p_A$ followed by $p_B$):
\begin{align*}
V_A &= p_A(V_0 + p_A) = p_A V_0 + p_A^2 \\
V_{AB} &= p_B(V_A + p_B) = p_B(p_A V_0 + p_A^2 + p_B) = p_A p_B V_0 + p_A^2 p_B + p_B^2
\end{align*}
Evaluate sequence B ($p_B$ followed by $p_A$):
\begin{align*}
V_B &= p_B(V_0 + p_B) = p_B V_0 + p_B^2 \\
V_{BA} &= p_A(V_B + p_A) = p_A(p_B V_0 + p_B^2 + p_A) = p_A p_B V_0 + p_A p_B^2 + p_A^2
\end{align*}
Compute the delta $\Delta = V_{AB} - V_{BA}$:
\begin{align*}
\Delta &= (p_A^2 p_B + p_B^2) - (p_A p_B^2 + p_A^2) \\
\Delta &= p_A p_B (p_A - p_B) - (p_A^2 - p_B^2) \\
\Delta &= p_A p_B (p_A - p_B) - (p_A - p_B)(p_A + p_B) \\
\Delta &= (p_A - p_B)(p_A p_B - (p_A + p_B))
\end{align*}
For $p_A, p_B \ge 2$ and $p_A \neq p_B$, $p_A p_B > p_A + p_B$, meaning the second term is strictly positive. Since $p_A - p_B \neq 0$, the delta $\Delta \neq 0$. 

Therefore, $V_{AB} \neq V_{BA}$. The PARM framework mathematically enforces covenantal sequence; the order of events fundamentally alters the ontological structure of the outcome.
\end{proof}

\subsection*{A.4 Well-Definedness of the Resonance Quotient ($R_Q$)}

The Multiplicity Resonance Quotient is defined as:
\[
R_Q = \frac{V_{\text{actual}} - V_{\min}}{V_{\max} - V_{\min}}
\]
For $R_Q$ to form a valid metric space across all permutations, we must prove that the denominator is strictly non-zero for any sequence containing at least two distinct letters.

\begin{lemma}[Strict Variance of Permutation Bounds]
For any sequence $W$ containing at least two distinct primes, $V_{\max} > V_{\min}$. 
\end{lemma}

\begin{proof}
Assume, for contradiction, that $V_{\max} = V_{\min}$. This implies that for all permutations $Q \in S(W)$, $V(Q)$ yields the exact same value. 
However, Theorem A.3 proves that swapping any two adjacent distinct letters in the flow alters the accumulated value by $\Delta = (p_A - p_B)(p_A p_B - (p_A + p_B)) \neq 0$. 
Furthermore, altering the final letter (the Seal) applies a multiplier of $p_N^2$ instead of $p_N$, massively scaling the outcome. Since $S(W)$ contains at least two distinct permutations (due to the presence of at least two distinct primes), their sealed states cannot be equal.
Thus, $V_{\max} \neq V_{\min}$, meaning $V_{\max} - V_{\min} > 0$. The denominator is strictly positive, and $R_Q \in [0,1]$ is well-defined.
\end{proof}

\section*{References}
\addcontentsline{toc}{section}{References}

\begin{thebibliography}{99}

\bibitem{bosworth2025anchor}
Bosworth, K. (2025).
\textit{Prime as Anchor: Beyond Detection to Deterministic Emergence}.
PhilArchive preprint. \url{https://philarchive.org/archive/BOSPAA-6}

\bibitem{bosworth2025threshold}
Bosworth, K. (2025).
\textit{The Threshold of Recursion: Why PAS > 0}.
PhilArchive preprint. \url{https://philarchive.org/archive/BOSTTO-9v1}

\bibitem{deleuze1969multiplicities}
Deleuze, G. (1969).
\textit{Multiplicities (Lecture 01, 30 November 1969)}.
Purdue University, Deleuze Seminars. \url{https://deleuze.cla.purdue.edu/lecture/lecture-01-11/}

\bibitem{laitman2016mathematical}
Laitman, M. (2016).
\textit{Hebrew – A Mathematical Language}.
Laitman.com. \url{https://laitman.com/2016/12/hebrew-a-mathematical-language/}

\bibitem{multiplicity2026parm}
Multiplicity Theory Foundation. (2026).
\textit{Execution Artifacts: Position-Aware Recursive Multiplicity (PARM) Core Operations}.
Archived code execution outputs (corruption\_test.csv, semantic\_rq\_sweep.csv).

\bibitem{murty2023pushdown}
Murty, S., et al. (2023).
\textit{Pushdown Layers: Encoding Recursive Structure in Transformer Language Models}.
Proceedings of EMNLP 2023. \url{https://github.com/MurtyShikhar/Pushdown-Layers}

\bibitem{savant2025prime}
Gelder, R. V. (2025).
\textit{Prime Indexing: The Simplest Filter for Lawful Recursion}.
LinkedIn Academic Publications. \url{https://www.linkedin.com/posts/multiplicity_multiplicity-primeindexing-recursion-activity-7374825373362077696-bYRV}

\bibitem{savant2025tensor}
Gelder, R. V. (2025).
\textit{Introducing Prime-Indexed Recursive Tensor Mathematics for Self-Stabilizing Systems}.
Multiplicity Network Publications. \url{https://www.linkedin.com/posts/multiplicity_deepinvent4good-inventathon-activity-7364106803887894528-0d1A}

\bibitem{shore2014precision}
Shore, H. (2014).
\textit{The Mathematical Precision of Biblical Hebrew}.
Haim Shore Blog. \url{https://haimshore.blog/2014/04/28/mathematical-precision-of-biblical-hebrew/}

\bibitem{zhang2025rlm}
Zhang, A. (2025).
\textit{RLM: Recursive Language Models}.
GitHub Repository. \url{https://github.com/alexzhang13/rlm}

\end{thebibliography}

\end{document}